\section{Appendix: SDPO-Fail, MOPD-Succeed Case Studies}
\label{sec:appendix-case-studies}

We focus here on prompts where the SDPO checkpoint still fails but the MOPD
checkpoint already succeeds.
Rather than emphasizing later rollout repair, these examples isolate a cleaner
checkpoint-level qualitative gap: both models see the same task, but MOPD is
more likely to preserve the crucial semantic distinction needed for a correct
program.

\subsection{Case Study 1: MOPD succeeds on a prompt where SDPO keeps the identifier-space error}

\paragraph{Candidate pair.}
The bib-mapping task provides a direct SDPO-versus-MOPD comparison.
At line 92, the SDPO checkpoint sample in
\texttt{diversoty\_sdpo\_50/50.jsonl} fails with accuracy $0.175$,
while the MOPD checkpoint sample in
\texttt{diversoty\_mopd\_50/50.jsonl} solves the same prompt perfectly.

\paragraph{SDPO failure.}
The SDPO trajectory mixes up the role of the permutation arrays and writes:
\begin{verbatim}
bib_to_person = [0] * (N + 1)
for i in range(N):
    bib_to_person[Q[i]] = P[i]

result = [str(bib_to_person[Q[i]]) for i in range(N)]
\end{verbatim}
This is incorrect in two ways.
First, \texttt{P[i]} is the stared-at person, not the wearer of bib
\texttt{Q[i]}; second, the final answer again stops in person-ID space instead
of projecting back to the bib number of the target person.
The resulting program preserves the rough shape of the mapping task but
composes the wrong relations.

\paragraph{MOPD success.}
The corresponding MOPD sample keeps the two identifier spaces separate:
\begin{verbatim}
bib_to_person = [0] * (N + 1)
for i in range(N):
    bib_to_person[Q[i]] = i + 1

for i in range(1, N + 1):
    person = bib_to_person[i]
    target_person = P[person - 1]
    result.append(str(Q[target_person - 1]))
\end{verbatim}
This trajectory performs the full chain correctly:
\texttt{bib} $\rightarrow$ \texttt{wearing person}
$\rightarrow$ \texttt{stared-at person}
$\rightarrow$ \texttt{target bib}.

\paragraph{Takeaway.}
This example shows a checkpoint-level separation rather than a later repair:
SDPO still exhibits identifier-space confusion, whereas MOPD already produces
the correct symbolic composition on the same prompt.

\subsection{Direct training-time evidence: the MOPD teacher is explicitly exposed to the same identifier-confusion failure}

\paragraph{A same-step mixed peer pool already exists at the start of training.}
The bib-mapping prompt above is not only a checkpoint-level contrast; it also
appears as a mixed success/failure peer group inside the MOPD rollout stream at
training step 1.
In
\texttt{.../PTrue-ATrue-FTrue-.../1.jsonl},
this prompt has eight peers in the same rollout group:
seven fail and one succeeds
(the successful sample is line 189; failed peers include lines 14, 15, 51, 78,
115, 179, and 243).
The failed peers are not generic failures.
Their reasoning repeatedly gets stuck on the same identifier-space ambiguity,
e.g., whether ``the person wearing bib $i$'' means person $i$ or the inverse
mapping $Q_j = i$.

\paragraph{The training log shows that MOPD is actually using failed peers as teacher context.}
At the corresponding MOPD training step, the console log records
\texttt{self\_distillation/reprompt\_sample\_fraction = 1.0} and
\texttt{self\_distillation/failure\_solution\_fraction = 0.99609375}
(\texttt{wandb/run-20260504\_205407-zwlvjx6t/files/output.log}, step 1).
This matters because the trainer code constructs a reprompt of the form
\begin{verbatim}
{prompt}{solution}{solution_summaries}{another_solution}{failure}{feedback}
\end{verbatim}
and, when failure solutions are enabled with
\texttt{failure\_solution\_condition = always}, inserts
\begin{verbatim}
A known failed solution (avoid repeating this mistake):
\end{verbatim}
followed by a failed peer trajectory
(\texttt{verl/trainer/ppo/ray\_trainer.py};
\texttt{verl/workers/config/actor.py}).
Taken together, the rollout file and the step-1 metrics imply that the
successful bib-mapping sample is trained under a teacher prompt that almost
certainly contains a same-task failed peer exhibiting the very same
identifier-confusion pattern.

\paragraph{Why this is the cleanest causal evidence we have.}
This is stronger than a mere before/after checkpoint comparison.
The failure pattern is already present in the peer pool at the same training
step, and the MOPD trainer is configured to expose such failed peers directly
in the teacher context.
We do not log per-token teacher logits in the current pipeline, so we cannot
print the exact top-$k$ distribution shift at the offending identifier tokens
from the training log alone.
But the causal path is explicit:
the teacher sees a concrete failed peer of the same bug family and supervises
the student on a successful target in that augmented context.

\paragraph{SDPO lacks this channel on the same training step.}
The matched SDPO run shows the same
\texttt{self\_distillation/reprompt\_sample\_fraction = 1.0} at step 1, but
\texttt{self\_distillation/failure\_solution\_fraction = 0.0}
(\texttt{wandb/run-20260505\_001730-iob9ijsm/files/output.log}, step 1).
This difference is exactly what the run configuration encodes:
the SDPO script defaults to
\texttt{INCLUDE\_FAILURE\_SOLUTION=False}, whereas the MOPD-style run enables
failed-peer context.
So for this prompt, MOPD has a direct training-time mechanism for saying
``here is a failed peer that confuses the identifier spaces; avoid this'',
while SDPO does not.

\subsection{Case Study 2: MOPD restores the dropped middle-character constraint}

\paragraph{Candidate pair.}
The evenly spaced \texttt{A}-\texttt{B}-\texttt{C} counting task gives a second
clean contrast.
At line 130, the SDPO checkpoint sample in
\texttt{diversoty\_sdpo\_50/50.jsonl} fails with accuracy $0.825$,
while the MOPD checkpoint sample in
\texttt{diversoty\_mopd\_50/50.jsonl} is fully correct.

\paragraph{SDPO failure.}
The SDPO code enforces equal spacing and checks the two endpoints, but drops the
requirement that the middle character must be \texttt{B}:
\begin{verbatim}
for j in range(1, n-1):
    for i in range(j):
        if S[i] == 'A':
            k = 2 * j - i
            if k < n and S[k] == 'C':
                count += 1
\end{verbatim}
This overcounts invalid triples such as \texttt{A-R-C} in the test case
\texttt{ARC}, where the correct answer is $0$.
The failure is therefore not an algorithmic collapse but a local semantic
omission: the code keeps the positional pattern and forgets one label
constraint from the specification.

\paragraph{MOPD success.}
The MOPD sample reinstates the missing check explicitly:
\begin{verbatim}
for j in range(1, n - 1):
    if s[j] == 'B':
        for i in range(j):
            if s[i] == 'A':
                k = 2 * j - i
                if k < n and s[k] == 'C':
                    count += 1
\end{verbatim}
The resulting program is still simple and structurally close to the SDPO
attempt, but it preserves the full symbolic constraint set.

\paragraph{Takeaway.}
This second example complements Case Study 1.
In Case Study 1, SDPO fails because it confuses representation spaces; here it
fails because it drops one clause from the natural-language specification.
MOPD succeeds in both settings, suggesting that its advantage is not tied to a
single failure type but extends across multiple forms of semantic fidelity.

\subsection{Case Study 3: The first repaired MOPD rollout appears early, while SDPO still fails late}

\paragraph{Candidate pair.}
The direction-opposite task gives a cleaner training-dynamics example.
At line 17, the SDPO checkpoint fails at both
\texttt{diversoty\_sdpo\_50/50.jsonl} and
\texttt{diversoty\_sdpo\_100/100.jsonl}, both with accuracy $0.0$.
The recorded failure mode also stays format-related:
the step-50 checkpoint is truncated, while the step-100 checkpoint is flagged
for missing a \texttt{```python} block.
In contrast, the MOPD checkpoint at
\texttt{diversoty\_mopd\_50/50.jsonl} already solves the same prompt with
accuracy $1.0$.

\paragraph{Early MOPD rollout repair.}
In the corresponding MOPD rollout stream, the first occurrence we find for this
prompt is already correct at rollout step 10
(\texttt{.../PTrue-ATrue-FTrue-.../10.jsonl}, line 1), with score $1.0$ and
accuracy $1.0$.
This makes step 10 the earliest observed point at which the repaired behavior is
visible in training-time rollouts for this prompt.

\paragraph{Late SDPO failure.}
The corresponding SDPO rollout family still shows the same bug much later.
At rollout step 80
(\texttt{.../PTrue-AFalse-FFalse-.../80.jsonl}, line 65), the model again fails
with accuracy $0.0$, and the environment feedback is
\texttt{Incorrect Format: Put your code inside a ```python ... ``` block.}

\paragraph{Takeaway.}
For this prompt, the qualitative split is not that MOPD invents a deeper
algorithm.
Instead, MOPD begins to internalize the basic answer-format discipline by about
step 10, while SDPO still has not stabilized that behavior even by the late
rollout/checkpoint regime.

\subsection{Case Study 4: MOPD shows an early correct rollout, whereas SDPO remains brittle late in training}

\paragraph{Candidate pair.}
The oyster-identification task shows the same pattern on a second prompt.
At line 98, both SDPO checkpoints again fail:
\texttt{diversoty\_sdpo\_50/50.jsonl} has accuracy $0.0$ with a truncated
attempt, and \texttt{diversoty\_sdpo\_100/100.jsonl} still has accuracy $0.0$
with an incorrect-format error.
The matching MOPD checkpoint in
\texttt{diversoty\_mopd\_50/50.jsonl} is fully correct.

\paragraph{Early MOPD signal.}
For this prompt, the earliest matching MOPD rollout we observe is already
correct at rollout step 5
(\texttt{.../PTrue-ATrue-FTrue-.../5.jsonl}, line 3), with score $1.0$ and
accuracy $1.0$.
This is the earliest observed rollout step where the repaired behavior appears
for this prompt.

\paragraph{Late SDPO failure.}
SDPO remains unstable on the same task much later in training.
At rollout step 55
(\texttt{.../PTrue-AFalse-FFalse-.../55.jsonl}, line 14), the model still
collapses into a fully truncated response, receiving score $0.0$ and accuracy
$0.0$.

\paragraph{Caveat and takeaway.}
This repair is not perfectly monotone across every rollout sample:
the MOPD stream also contains a later wrong-answer trace for the same prompt at
step 35.
Still, the earliest successful rollout appears very early on the MOPD side,
whereas SDPO continues to show the same format-control failure late in training
and still does not solve the checkpoint evaluation at steps 50 or 100.
